% =============================================================================
% Chapter 2 --- Canonical data schema
% =============================================================================
\chapter{Canonical data schema}
\label{ch:schema}

The structured layer uses four Draft-07 JSON Schemas and a central
\corpus{structured/enums.yaml} for controlled vocabulary.

\section{Schema inventory}

\begin{table}[h]
\centering
\small
\begin{tabularx}{\textwidth}{p{5.2cm}X}
\toprule
\textbf{Schema} & \textbf{Describes} \\
\midrule
\tool{regulation.schema.json}        & One published framework, standard, guideline, research paper, or index. Captures publisher, version, scope, section tree, and a hash of the source PDF. \\
\tool{requirement.schema.json}       & One atomic requirement or observational claim. Carries \texttt{implementation[]} and \texttt{verification[]} arrays that bind it to the repository and to the conformance tools. \\
\tool{conformance-tool.schema.json}  & One vendored testing framework. Enumerates plugins, benchmarks, and integration points. \\
\tool{corpus-index.schema.json}      & The top-level index that lists every item in the corpus, typed as regulation, requirement, conformance\_tool, or rag\_artifact. \\
\bottomrule
\end{tabularx}
\caption{Schemas under \corpus{structured/schema/}.}
\label{tab:schemas}
\end{table}

\section{Requirement record --- the linkage schema}

The requirement schema is the key extension that separates this corpus
from the earlier finance corpora. Each requirement carries:

\begin{itemize}
  \item \texttt{source}: regulation id, section, page, and (optionally) a
        verbatim quote.
  \item \texttt{normativeText}: a paraphrased imperative.
  \item \texttt{category} and \texttt{obligation}: structured priority
        (\emph{must}, \emph{should}, \emph{may}, \emph{descriptive}).
  \item \texttt{lifecyclePhase}: where in the agent lifecycle the
        requirement applies (planning, design, development,
        pre-deployment, deployment, post-deployment, decommission).
  \item \texttt{implementation[]}: an array of pointers to code, config,
        documentation, schema, or dataset evidence. A \texttt{kind: "none"}
        entry explicitly records a known gap.
  \item \texttt{verification[]}: an array of pointers to AI Verify plugins
        (by \texttt{plugin}), Moonshot benchmarks (by \texttt{testId}),
        repository tests (by \texttt{path}), manual-review entries, or
        process-checklist items.
  \item \texttt{status}: \emph{implemented}, \emph{partial}, \emph{gap},
        \emph{planned}, or \emph{not\_applicable}.
\end{itemize}

\begin{lstlisting}[language=jsonx,caption={Example requirement with
implementation and verification arrays (excerpt).},label={lst:req-example}]
{
  "requirementId": "REG-MGF-2.1.2-001",
  "source": {
    "regulationId": "mas-mgf-agentic-ai-v1",
    "section": "2.1.2",
    "page": 11,
    "quote": "Bound risks through design by defining agents' limits ..."
  },
  "normativeText": "The agent must operate against an explicit allow-list
                    of tools and external systems; broader access must be
                    denied by default.",
  "category": "bound_risks",
  "obligation": "must",
  "lifecyclePhase": "design",
  "implementation": [
    {"kind": "code",
     "path": "src/intelligence/ai-core-pal/mcp_server/btp_pal_mcp_server.py",
     "symbol": "btp_pal_mcp_server"},
    {"kind": "config",
     "path": "src/intelligence/ai-core-pal/data_products/registry.yaml"}
  ],
  "verification": [
    {"kind": "test",
     "path": "src/intelligence/ai-core-pal/tests/test_mcp_server.py"},
    {"kind": "moonshot-test", "testId": "adversarial_prompts"}
  ],
  "status": "partial",
  "priority": "critical"
}
\end{lstlisting}

\section{Conformance-tool record}

The conformance-tool schema captures the tools themselves so that
requirement records can reference them by stable identifiers. Every
\texttt{plugin.pluginId} and \texttt{benchmark.benchmarkId} is part of the
vocabulary that requirement records can cite.

\begin{lstlisting}[language=jsonx,caption={Excerpt of
\corpus{structured/conformance-tools/aiverify.json}.}]
{
  "toolId": "aiverify",
  "vendoredPath": "docs/regulations/aiverify-main",
  "kind": "testing-framework",
  "plugins": [
    { "pluginId": "aiverify.stock.process-checklist",
      "path": "docs/regulations/aiverify-main/stock-plugins/...",
      "purpose": "11-principle process-checklist questionnaire ...",
      "algorithmType": "process" },
    { "pluginId": "aiverify.stock.robustness-toolbox",
      "algorithmType": "robustness" }
    // ... nine more plugins
  ]
}
\end{lstlisting}

\section{Enum vocabularies}

\corpus{structured/enums.yaml} centralises the controlled vocabularies
that recur across schemas. The authoritative values live in the schemas
themselves; \texttt{enums.yaml} is the human index.

\begin{table}[h]
\centering
\small
\begin{tabularx}{\textwidth}{lX}
\toprule
\textbf{Enum} & \textbf{Values} \\
\midrule
Regulation type & framework, standard, guideline, statute, research, index, toolkit \\
Normativity    & mandatory, recommended, advisory, descriptive \\
Obligation verb & must, should, may, descriptive \\
Requirement category & bound\_risks, human\_accountability, technical\_controls, end\_user\_responsibility, transparency, evaluation, observation \\
Lifecycle phase & planning, design, development, pre\_deployment, deployment, post\_deployment, decommission \\
Implementation kind & code, config, doc, schema, dataset, none \\
Verification kind & aiverify-plugin, moonshot-test, test, manual-review, process-checklist, external-audit \\
Status & implemented, partial, gap, planned, not\_applicable \\
\bottomrule
\end{tabularx}
\caption{Controlled vocabularies.}
\label{tab:enums}
\end{table}
\clearpage
